\documentclass[11pt]{amsart}

\usepackage[T1]{fontenc}
\usepackage{lmodern}
\usepackage{microtype}
\usepackage{amsmath,amssymb,amsthm}
\usepackage{booktabs}
\usepackage[hidelinks]{hyperref}

\hypersetup{
  pdftitle={A Four-Vertex Counterexample to Conjecture 7.6 of Devriendt}
}

\newtheorem{theorem}{Theorem}
\newtheorem{proposition}[theorem]{Proposition}

\newcommand{\one}{\mathbf{1}}

\title[A four-vertex counterexample]
{A Four-Vertex Counterexample to Conjecture 7.6 of Devriendt}
\date{}

\begin{document}

\begin{abstract}
Conjecture~7.6 of Devriendt asserts that a normalized weighted graph with
submodular resistance capacity has nonnegative resistance curvature. We
give an exact counterexample on the diamond graph $K_4-e$: the weighting is
normalized and its resistance capacity is submodular, but its curvature
vector is $(2/5,2/5,7/30,-1/30)$. The example is vertex-minimal. It refutes
Conjecture~7.6 but does not refute the graph-level Conjecture~7.7.
\end{abstract}

\maketitle

\section{The counterexample}

Let $G=(V,E)$ be a connected simple graph with positive conductances $c$.
Write $\Omega=(\omega_{uv})_{u,v\in V}$ for its effective-resistance matrix
and $K=\Omega^{-1}$. Following \cite{Devriendt}, the weighting is
\emph{normalized} if
\[
\one^{\mathsf T}K\one=1.
\]
For a normalized weighting, the resistance capacity is the set function
$\tau:2^V\to\mathbb R$ defined by
\[
\tau_{\varnothing}=0,
\qquad
\tau_{\{v\}}=\frac12,
\qquad
\tau_U=\frac12+\frac12
\left(\one^{\mathsf T}\Omega[U]^{-1}\one\right)^{-1}
\quad (|U|\ge2),
\]
where $\Omega[U]$ is the principal submatrix indexed by $U$. The resistance
curvature is
\[
p_v=1-\frac12\sum_{e\ni v}c_e\omega_e.
\]
Conjecture~7.6 of \cite{Devriendt} states that, for a normalized weighting,
submodularity of $\tau$ implies $p\ge0$.

\begin{theorem}
Let $V=\{0,1,2,3\}$ and
\[
E=\{02,03,12,13,23\},
\]
so that $G$ is the diamond graph with nonedge $01$. In the edge order
$(02,03,12,13,23)$, set
\[
c=\frac{109}{300}(1,2,1,2,2).
\]
Then $c$ is normalized and $\tau$ is submodular, whereas
\[
p=\left(\frac25,\frac25,\frac7{30},-\frac1{30}\right).
\]
Consequently, Conjecture~7.6 is false.
\end{theorem}

\begin{proof}
The weighted Laplacian is
\[
L=\frac{109}{300}
\begin{pmatrix}
 3& 0&-1&-2\\
 0& 3&-1&-2\\
-1&-1& 4&-2\\
-2&-2&-2& 6
\end{pmatrix}.
\]
Using
\[
L^{\dagger}
=
\left(L+\frac14\one\one^{\mathsf T}\right)^{-1}
-\frac14\one\one^{\mathsf T}
\]
and
$\omega_{uv}=(e_u-e_v)^{\mathsf T}L^{\dagger}(e_u-e_v)$ gives
\[
\Omega=\frac1{109}
\begin{pmatrix}
  0&200&140&110\\
200&  0&140&110\\
140&140&  0& 90\\
110&110& 90&  0
\end{pmatrix}.
\]
Direct inversion yields
\[
K\one=
\left(\frac25,\frac25,\frac7{30},-\frac1{30}\right)^{\mathsf T},
\qquad
\one^{\mathsf T}K\one=1,
\]
so the weighting is normalized. Moreover, in the stated edge order,
\[
(c_e\omega_e)_{e\in E}
=
\frac1{15}(7,11,7,11,9).
\]
Substitution in the definition of curvature gives the displayed vector
$p$.

It remains to verify submodularity. For distinct $i,j\in V$ and
$S\subseteq V\setminus\{i,j\}$, define the elementary slack
\[
\Delta_{ij}(S)
=
\tau_{S\cup\{i\}}+\tau_{S\cup\{j\}}
-\tau_S-\tau_{S\cup\{i,j\}}.
\]
A set function is submodular if and only if all such slacks are
nonnegative. Put $\Gamma_{ij}(S)=3924\,\Delta_{ij}(S)$. Substitution of
$\Omega$ into the definition of $\tau$ gives the following complete list
up to the automorphism $0\leftrightarrow1$:

\begin{center}
\small
\renewcommand{\arraystretch}{1.12}
\begin{tabular}{ccrrrr}
\toprule
$ij$ & $\{a,b\}=V\setminus\{i,j\}$
& $\Gamma_{ij}(\varnothing)$
& $\Gamma_{ij}(\{a\})$
& $\Gamma_{ij}(\{b\})$
& $\Gamma_{ij}(\{a,b\})$\\
\midrule
$01$ & $\{2,3\}$ & 162  & 560  & 165 & 0\\
$02$ & $\{1,3\}$ & 702  & 1100 & 414 & 249\\
$03$ & $\{1,2\}$ & 972  & 975  & 684 & 124\\
$23$ & $\{0,1\}$ & 1152 & 864  & 864 & 13\\
\bottomrule
\end{tabular}
\end{center}

The rows for $12$ and $13$ are obtained from those for $02$ and $03$ by
swapping $0$ and $1$. Thus the table accounts for all $24$ elementary
slacks, each of which is nonnegative. Hence $\tau$ is submodular, while
$p_3=-1/30<0$.
\end{proof}

\begin{proposition}
No positively weighted connected simple graph on two or three vertices has
negative resistance curvature. Hence four is the smallest possible order
of a counterexample.
\end{proposition}

\begin{proof}
For an edge $e$, the quantity $c_e\omega_e$ is its inclusion probability
in the weighted random spanning tree. It is therefore $1$ on a bridge and
strictly less than $1$ on a nonbridge. The connected simple graphs on two
or three vertices are $K_2$, $P_3$, and $K_3$. The first two are trees and
have curvature vectors $(1/2,1/2)$ and $(1/2,0,1/2)$, respectively. In
$K_3$, each vertex is incident to two nonbridges, so the sum of the two
incident quantities $c_e\omega_e$ is strictly less than $2$ and every
curvature coordinate is positive.
\end{proof}

\medskip
\noindent\textbf{Scope.}
Conjecture~7.7 of \cite{Devriendt} asks whether a graph admits a normalized
weighting with submodular resistance capacity exactly when it admits a
weighting with nonnegative resistance curvature. It does not require the
same weighting in the two existence statements. The diamond graph is
resistance positive: unit conductances give curvature
\[
\left(\frac38,\frac38,\frac18,\frac18\right),
\]
and a common rescaling normalizes the weighting without changing curvature.
Thus the example refutes precisely the fixed-weight implication in
Conjecture~7.6; it does not refute Conjecture~7.7.

\begin{thebibliography}{1}

\bibitem{Devriendt}
K.~Devriendt,
\emph{Graphs with nonnegative resistance curvature},
Ann. Comb. \textbf{30} (2026), 415--438.
\href{https://doi.org/10.1007/s00026-025-00774-x}
{doi:10.1007/s00026-025-00774-x}.

\end{thebibliography}

\end{document}